\chapter{Fazit und Ausblick}
\label{chap:fazit}

\section{Zusammenfassung und Schlussfolgerung}
\label{sec:fazit_zusammenfassung}

Die vorliegende Arbeit hat die Erkennung von DDoS-ähnlichen Angriffsmustern als sequenzielles Entscheidungsproblem in einer dynamischen, verrauschten Netzwerkumgebung modelliert und mit einem Double Deep Q-Network (DDQN) umgesetzt. Ausgangspunkt war die in Kapitel~\ref{chap:einleitung} formulierte Fragestellung, inwieweit KI-basierte Verfahren unter realitätsnahen Störbedingungen eine belastbare Angriffserkennung ermöglichen. Auf Basis der in Kapitel~\ref{chap:grundlagen} und Kapitel~\ref{chap:related_work} diskutierten methodischen Grundlagen wurde in Kapitel~\ref{chap:datensatz} eine konsistente Datenpipeline aufgebaut und in Kapitel~\ref{chap:implementierung} eine reproduzierbare RL-Umgebung mit AWGN-Modellierung implementiert. Kapitel~\ref{chap:evaluierung} hat das Evaluationsprotokoll und die Metriken festgelegt, Kapitel~\ref{chap:ergebnisse} die Leistungsfähigkeit des Ansatzes im Vergleich mit den Baselines empirisch quantifiziert.

Die Ergebnisse zeigen, dass der Double-DQN-Agent (DDQN) bei 30\,dB einen F1-Score von $0{,}8935$ erreicht (Precision $0{,}8853$, Recall $0{,}9018$). Gegenüber Random Forest, SVM und MLP liegt er auf allen vier SNR-Stufen vorn. Der Vorsprung gegenüber dem MLP beträgt allerdings nur $0{,}25$ bis $0{,}62$ Prozentpunkte und bewegt sich in der Größenordnung der Seed-Streuung. Belastbar ist die Rangfolge DDQN $\gtrsim$ MLP $>$ SVM $>$ Random Forest, nicht ein signifikanter Vorsprung. Die Equalized Loss of Accuracy nach Sáez et al. bestätigt diese Rangfolge auf allen SNR-Stufen \autocite{saezEvaluatingClassifierBehavior2016}. Bei 5\,dB verlieren alle vier Verfahren zwischen $18{,}8$ und $20{,}2\,\%$ ihres F1-Scores. Der Verlust folgt damit überwiegend aus der Klassenüberlappung des Kanals, nicht aus dem Modell. Dass der Random Forest hier zurückliegt, während er in der Jamming-Erkennung von Arjoune et al. vorn liegt \autocite{arjouneNovelJammingAttacks2020}, erklärt sich aus der nahezu linearen, schräg im Merkmalsraum liegenden Bayes-Grenze des Generatormodells (Abschnitt~\ref{sec:ergebnisse_literatur}).

Über die untersuchten SNR-Stufen von 30\,dB bis 5\,dB fällt der F1-Score monoton ab, zunächst moderat von $0{,}8935$ auf $0{,}8298$ bei 10\,dB, bei 5\,dB dann auf $0{,}7260$. Der Agent bleibt damit bis etwa 10\,dB operationell nutzbar. Der Recall liegt auf jeder Stufe über der Precision, weil die asymmetrische Reward-Matrix ($R_{FN}=-2$, $R_{FP}=-1$) die Entscheidungsschwelle auf eine Angriffs-Posterior von $1/3$ senkt. Bei sinkendem SNR bricht deshalb die Precision stärker ein als der Recall. Für den Netzbetrieb heißt das: Unter schlechten Kanalbedingungen steigt die Zahl der Fehlalarme schneller als die Zahl verpasster Angriffe.

Im Gesamtfazit bestätigt die Arbeit die Eignung des RL-basierten Ansatzes für die adaptive DDoS-Erkennung in dynamischen Kommunikationsnetzen. Die Kombination aus sequenzieller Entscheidungslogik, Experience Replay, Zielnetzwerk und rauschbehafteter Umweltsimulation liefert einen methodischen Mehrwert gegenüber rein statischen Klassifikationsverfahren. Zugleich bleiben Grenzen sichtbar: Die Aussagekraft ist an das gewählte Simulationsdesign, die betrachteten Merkmale und die Abdeckung der Datensätze gebunden. Für eine weitergehende Generalisierung sind daher zusätzliche Untersuchungen auf heterogeneren Lastprofilen, realen Traffic-Spuren und erweiterten Angriffstypen notwendig \autocite{balaAITechniquesIoTbased2024, kumarComprehensiveReviewVulnerabilities2024}.

\section{Ausblick und Forschungsperspektiven}
\label{sec:fazit_ausblick}

Aus den gewonnenen Befunden ergeben sich mehrere fundierte Forschungsperspektiven für zukünftige Arbeiten:

\begin{enumerate}
    \item \textbf{Validierung auf den Originaldatensätzen:} Die Experimente dieser Arbeit nutzen ein parametrisches Flussmodell, dessen Merkmalsstruktur sich an DDoS-AT-2022 \autocite{mittalDDoSAT2022DistributedDenial2023} und IoT-DH \autocite{saifIoTDHDatasetClassification2024} orientiert (Abschnitt~\ref{sec:generatormodell}). Der nächste Schritt ist die Wiederholung von Training und Test auf den Originaldaten beider Datensätze, mit einem getrennten Validierungs-Split, einem AWGN-Kanal auf den realen Merkmalen und denselben Baselines. Erst dieser Vergleich zeigt, ob die Rangfolge und die SNR-Robustheit auf realem Verkehr bestehen bleiben. Dabei lässt sich zugleich prüfen, ob die Reward-getriebene Verschiebung zum Recall auch bei den realen Klassenverhältnissen (z.\,B. rund $37\,\%$ Angriffsanteil in IoT-DH) erhalten bleibt.
    \item \textbf{Erweiterung auf komplexe Netzwerktopologien:} Der Ansatz sollte auf komplexere Topologien mit hierarchischen Segmenten, variabler Linkqualität und verteilten Kontrollinstanzen übertragen werden, um die Skalierbarkeit in realen Betreiberumgebungen zu validieren.
    \item \textbf{Einsatz fortgeschrittener RL-Architekturen:} Zur Steigerung der Lernstabilität und Sample-Effizienz bietet sich \textit{Prioritized Experience Replay} an. Sinnvoll wird der Bootstrapping-Anteil des DDQN erst bei mehrschrittigen Formulierungen, etwa wenn der Agent Verkehrsfenster sequenziell beobachtet und Gegenmaßnahmen steuert.
    \item \textbf{Multi-Agenten-Systeme:} Die Untersuchung kooperativer Multi-Agenten-Architekturen, bei denen spezialisierte Agenten verteilt für Erkennung, Priorisierung und automatisierte Gegenmaßnahmen zusammenwirken, verspricht eine weitere Reduktion der Reaktionszeiten.
    \item \textbf{Gezielte Störungen und robustes Training:} Die gemessene Robustheit gilt nur für stochastisches AWGN. Der nächste Test ist ein Angriff mit FGSM und BIM auf denselben Agenten, wie ihn Merzouk et al. für einen DRL-Detektor durchgeführt haben \autocite{merzoukEvadingDeepReinforcement2022}. Als Gegenmaßnahme bietet sich der Regularisierer aus dem SA-MDP-Rahmen von Zhang et al. an, der die Änderung der gewählten Aktion innerhalb einer $\varepsilon$-Umgebung der Beobachtung begrenzt \autocite{zhangRobustDeepReinforcement2020}. Ran et al. setzen eine gegen gestörte Zustandsbeobachtungen robust trainierte DQN bereits zur Erkennung von False-Data-Injection-Angriffen in Stromnetzen ein \autocite{ranRobustDataDrivenAdversarial2024}. Ob sich der Ansatz auf Flussmerkmale übertragen lässt, ist offen.
    \item \textbf{Betriebskennzahlen und Erklärbarkeit:} Die Leistungsanalyse sollte um Latenz, Ressourcenverbrauch und Verfahren der \textit{Explainable AI} (XAI) ergänzt werden.
\end{enumerate}

In der Gesamtschau liefert die Arbeit damit einen konsistenten Beitrag zur KI-gestützten Angriffserkennung: Sie verbindet theoretische Fundierung, reproduzierbare Implementierung und empirische Robustheitsanalyse zu einem geschlossenen methodischen Rahmen. Die Ergebnisse unterstützen die zentrale Arbeitshypothese, dass DDQN-basierte Verfahren bei geeigneter Modellierung und Evaluation eine leistungsfähige und zugleich ausbaufähige Grundlage für adaptive DDoS-Abwehr in modernen Kommunikationsnetzen darstellen.